% Reconstructed text from V5 PDF page 28; non-native source
\par\smallskip\noindent\rule{\textwidth}{0.4pt}\par{\sffamily\footnotesize Original V5 page 28 -- recovered text}\par\begin{flushleft}\footnotesize\ttfamily
\par\smallskip
\par\smallskip
\noindent Diagnostics and Serving Infrastructure for MAS 28\par
\par\smallskip
\par\smallskip
\par\smallskip
\noindent  TokenCake [4] implements temporal-spatial dual scheduling for multi-agent KV cache man-\par
\noindent agement, coordinating two schedulers that operate on complementary dimensions. The tempo-\par
\noindent ral scheduler is event-driven: triggered by function-call start events from the frontend API, it\par
\noindent proactively offloads stalled agents’ KV caches to CPU when three conditions are met simulta-\par
\noindent neously: the predicted stall exceeds the round-trip transfer time, a waiting request can use the\par
\noindent freed blocks, and GPU pressure exceeds 60\% (§4). A composite soft score integrating GPU\par
\noindent pressure, block-fit quality, and upload safety margin governs offload priority, with critical-path\par
\noindent agents penalized to avoid offloading bottleneck agents. Predictive upload begins restoring KV\par
\noindent cache before the predicted function-call completion using dynamic EWMA forecasting (Equa-\par
\noindent tion 1).\par
\noindent  The spatial scheduler dynamically partitions GPU memory into a shared pool (all agents)\par
\noindent and a reserved pool (critical agents only), with watermark-driven adjustment: when GPU us-\par
\noindent age exceeds 75\%, the reserved ratio ρ increases by 0.05; when below 40\%, it decreases by\par
\noindent 0.05 (clamped to [0.05, 0.30]). A hybrid priority metric (Equations 5–6) combines structural\par
\noindent importance (DAG depth and degree), synchronization pressure (straggler boost at join points),\par
\noindent temporal aging, recomputation cost, and graph context.\par
\noindent  The system achieves up to 47.06\% end-to-end latency reduction vs. vLLM in memory-\par
\noindent constrained settings (gpu\_mem\_utilization=0.5, §7.2, Figure 9) and a 16.9 percentage-point\par
\noindent GPU utilization improvement (85.7–87.0\% vs. 69.9–74.1\% for vLLM). Round-trip KV migra-\par
\noindent tion takes 15.8–79.8 ms compared to 593–2,300 ms for recomputation (26.8–37.5× speedup,\par
\noindent §7.6, Figure 17). A non-monotonic sensitivity result is notable: intermediate prediction noise\par
\noindent (scale 0.25) causes an 8.3\% performance regression because mistimed uploads pass feasibility\par
\noindent checks but still cause harmful migrations (§7.5). This two-dimensional scheduling represents\par
\noindent the most granular eviction policy among surveyed systems (EuroSys 2027).\par
\noindent  CacheBlend [41] addresses the RAG serving scenario: pre-computed KV caches exist for\par
\noindent knowledge chunks, but naively concatenating chunks produces attention errors because each\par
\noindent chunk was computed in isolation without cross-chunk attention. CacheBlend selectively re-\par
\noindent computes a subset of “critical” tokens at each layer — those whose KV representations differ\par
\noindent most between isolated and concatenated contexts — while reusing the remaining cached val-\par
\noindent ues. This selective recomputation achieves 2.2–3.3× TTFT reduction and 2.8–5× throughput\par
\noindent improvement versus full recomputation, with no quality degradation on open-source LLMs (Eu-\par
\noindent roSys 2025).\par
\noindent  TokenDance enables collective KV cache sharing across synchronized multi-agent rounds.\par
\noindent Rather than storing independent KV caches per agent, it maintains a single master copy plus\par
\noindent block-sparse diffs per agent — a diff-aware storage scheme. When agents share similar contexts\par
\noindent (common in round-based simulations where agents observe the same environment), the shared\par
\noindent blocks are stored once regardless of agent count. This achieves 2.7× more concurrent agents\par
\noindent versus vLLM with prefix caching, 17.5× per-agent KV storage reduction, and 1.9× prefill\par
\noindent speedup. Diff compression ratios of 11–17× are reported on GenerativeAgents and AgentSoci-\par
\noindent ety benchmarks. The approach is tailored to the All-Gather communication pattern and requires\par
\noindent a central scheduler; arbitrary agent topologies are not supported.\par
\noindent  ForkKV [36] targets the multi-LoRA agent serving scenario, where multiple agents share the\par
\noindent same base model but differ in their LoRA adapters. It splits KV cache into a shared base com-\par
\noindent ponent (from the pretrained model) and lightweight agent-specific residuals (LoRA-induced\par
\noindent divergences), indexed by a DualRadixTree. The copy-on-write deduplication exploits LoRA’s\par
\noindent low-rank structure: since LoRA adapters modify only a small subspace of the attention com-\par
\noindent putation, most KV cache entries are identical across adapters. ForkKV achieves up to 3.0×\par
\noindent throughput versus state-of-the-art multi-LoRA serving with negligible quality impact. Built on\par
\end{flushleft}
